\begin{proof}[Proof of \cref{obj:lem:one-cell-testing-family}]
\begin{enumerate}
\item \textbf{Construction of the family.}
Choose
\[
\underline c:=\frac1{256},
\qquad
\underline c_\alpha:=\frac{3(1-\alpha)^2}{512}
\quad(\alpha\in\mathbb R).
\]
These constants have the required positivity. Fix admissible \(n,d,q\), and write \(N=nq>0\). Define
\[
\operatorname{clip}_{[0,1]}(u):=\max\{0,\min\{1,u\}\}
\quad(u\in\mathbb R).
\]
For every \(u\in\mathbb R\), define \(P_u\) by
\[
X=0,\qquad S(0)=S(1)=Y(0)=0,
\]
\[
A\sim\operatorname{Bernoulli}(1/2),\qquad
Y(1)\sim\operatorname{Bernoulli}
       \bigl(\operatorname{clip}_{[0,1]}(u)\bigr),\qquad
R\sim\operatorname{Bernoulli}(q),
\]
with \(A,Y(1),R\) mutually independent. The first baseline category exists because \(d\ge1\). All atom weights are nonnegative, and their sum is
\[
\left(\frac12+\frac12\right)
\left(\operatorname{clip}_{[0,1]}(u)
      +1-\operatorname{clip}_{[0,1]}(u)\right)
\left(q+1-q\right)=1.
\]
Thus \(P_u\) is a full-data probability law for every real \(u\).

For \(u\in[1/4,3/4]\), clipping leaves \(u\) unchanged. Set the realized variables to \(S=S(A)\) and \(Y=Y(A)\), and use independent observed records with their induced marginal law. This gives the sampling and consistency conditions. Independence and the fair treatment coin give
\[
A\perp\!\!\!\perp
  \{X,S(0),S(1),Y(0),Y(1)\},
\qquad P_u(A=1)=\frac12.
\]
Moreover, for \(a,y,r\in\{0,1\}\),
\[
P_u(R=r,Y=y\mid X=0,A=a,S=0)
=
P_u(R=r)\,P_u(Y=y\mid A=a),
\]
because \(R\) is independent of \((A,Y(1))\) and \(Y=AY(1)\). This verifies \cref{obj:ass:arrival-mar} on both occupied cells; all joint probabilities on the other cells vanish. Direct summation also gives
\[
p_{axs}(P_u)=
\begin{cases}
1/2,&x=0,\ s=0,\\
0,&\text{otherwise},
\end{cases}
\qquad
P_u(A=a,X=x,S=s,R=1)=q\,p_{axs}(P_u).
\]
Consequently, \(\rho_{a00}(P_u)=q\), so
\cref{obj:ass:occupied-cell-arrival} holds. The fixed parameters satisfy
\cref{obj:ass:rare-arrival-slice}; hence
\[
P_u\in\mathcal M_{n,d,q}.
\]
Finally, the construction yields
\[
P_u\!\left(X=0,S(0)=S(1)=0,Y(0)=0\right)=1,
\qquad P_u(R=1)=q,
\qquad
\tau(P_u)=E_{P_u}[Y(1)-Y(0)]=u.
\]
The family is chosen independently of \(\alpha\).
% lean: CausalSmith.Stat.MarRareqLogfrontier.oneCellFamilyLaw_model
% lean: CausalSmith.Stat.MarRareqLogfrontier.oneCellFamilyLaw_support
% lean: CausalSmith.Stat.MarRareqLogfrontier.oneCellFamilyLaw_arrival_mass
% lean: CausalSmith.Stat.MarRareqLogfrontier.oneCellFamilyLaw_ate
% lean: CausalSmith.Stat.MarRareqLogfrontier.one_cell_testing_family

\item \textbf{Observed laws and chi-square divergence.}
Define the observed-record space, its one-record laws, and its sample laws by
\[
\mathcal O_d:=\{0,\ldots,d-1\}\times\{0,1\}^4,
\qquad
Q_u:=\mathcal L_{P_u}(X,A,S,R,RY),
\qquad
\mathbb Q_u:=Q_u^{\otimes n}.
\]
For \(u\in[1/4,3/4]\), the potentially nonzero atom probabilities are
\[
\begin{array}{c|c}
o=(X,A,S,R,RY)&Q_u(\{o\})\\ \hline
(0,0,0,0,0)&(1-q)/2\\
(0,1,0,0,0)&(1-q)/2\\
(0,0,0,1,0)&q/2\\
(0,1,0,1,0)&q(1-u)/2\\
(0,1,0,1,1)&qu/2
\end{array}
\]
and all remaining atoms have probability zero. Each full-data atom of positive \(P_u\)-mass has positive \(P_{1/2}\)-mass: the treatment weights are identical, both central outcome weights are positive, and a zero arrival weight vanishes under both laws. Passing to the observed marginal and then to products gives
\[
Q_u\ll Q_{1/2},
\qquad
\mathbb Q_u\ll\mathbb Q_{1/2}.
\]

For these finite laws, define
\[
\chi^2(Q_u\Vert Q_{1/2})
:=
\sum_{\substack{o\in\mathcal O_d\\Q_{1/2}(\{o\})>0}}
\frac{\bigl(Q_u(\{o\})-Q_{1/2}(\{o\})\bigr)^2}
     {Q_{1/2}(\{o\})}.
\]
Expanding the square and using the unit total masses gives
\[
1+\chi^2(Q_u\Vert Q_{1/2})
=
\sum_{\substack{o\in\mathcal O_d\\Q_{1/2}(\{o\})>0}}
\frac{Q_u(\{o\})^2}{Q_{1/2}(\{o\})}.
\]
If \(0<q<1\), summing the five displayed atoms yields
\[
1+\chi^2(Q_u\Vert Q_{1/2})
=
(1-q)+\frac q2+q(1-u)^2+qu^2
=
1+2q\left(u-\frac12\right)^2.
\]
If \(q=1\), the two nonarrival atoms have zero mass and are omitted; the remaining three terms give the same identity. Thus
\[
\chi^2(Q_u\Vert Q_{1/2})
=
2q\left(u-\frac12\right)^2.
\]

To record the product calculation, define
\[
\ell_u(o):=
\begin{cases}
Q_u(\{o\})/Q_{1/2}(\{o\}),&Q_{1/2}(\{o\})>0,\\
0,&Q_{1/2}(\{o\})=0,
\end{cases}
\qquad
L_{u,n}(\boldsymbol o):=\prod_{i=1}^{n}\ell_u(o_i)
\quad(\boldsymbol o\in\mathcal O_d^n).
\]
These are likelihood ratios almost surely under the corresponding reference laws. Independence gives
\[
E_{\mathbb Q_{1/2}}L_{u,n}=1,
\qquad
E_{\mathbb Q_{1/2}}L_{u,n}^2
=
\left(E_{Q_{1/2}}\ell_u^2\right)^n.
\]
Defining the sample chi-square divergence by
\[
\chi^2(\mathbb Q_u\Vert\mathbb Q_{1/2})
:=
E_{\mathbb Q_{1/2}}(L_{u,n}-1)^2,
\]
we obtain
\[
1+\chi^2(\mathbb Q_u\Vert\mathbb Q_{1/2})
=
\left(1+\chi^2(Q_u\Vert Q_{1/2})\right)^n
\le
\exp\!\left(n\chi^2(Q_u\Vert Q_{1/2})\right).
\]
The final inequality follows from \(1+x\le e^x\) for \(x\ge0\). With the total-variation convention of
\cref{obj:lem:randomized-kernel-tv-comparison}, the finite-space formula and Cauchy--Schwarz give
\[
\operatorname{TV}(\mathbb Q_u,\mathbb Q_{1/2})
=
\frac12E_{\mathbb Q_{1/2}}|L_{u,n}-1|
\le
\frac12
\sqrt{\chi^2(\mathbb Q_u\Vert\mathbb Q_{1/2})}.
\]
% lean: CausalSmith.Stat.MarRareqLogfrontier.oneCellFamilyLaw_obs_ac
% lean: CausalSmith.Stat.MarRareqLogfrontier.oneCellFamilyLaw_obs_atom
% lean: CausalSmith.Stat.MarRareqLogfrontier.oneCellFamilyLaw_obs_chi
% lean: Causalean.Stat.one_add_chiSqDiv_pi_iid_general
% lean: Causalean.Stat.tvDist_le_half_sqrt_chiSqDiv

\item \textbf{A two-point testing bound.}
Define
\[
\Delta_{\mathrm{pt}}
:=
\min\left\{\frac14,\frac1{4\sqrt N}\right\},
\qquad
u_0^{\mathrm{pt}}:=\frac12,
\qquad
u_1^{\mathrm{pt}}:=\frac12+\Delta_{\mathrm{pt}}.
\]
Since \(N>0\),
\[
0<\Delta_{\mathrm{pt}}\le\frac14.
\]
The two means are therefore distinct and belong to \([1/4,3/4]\). Also,
\[
\Delta_{\mathrm{pt}}^2\le\frac1{16N},
\qquad
\chi^2(Q_{u_1^{\mathrm{pt}}}\Vert Q_{1/2})
=
2q\Delta_{\mathrm{pt}}^2
\le\frac1{8n}.
\]
The product calculation above yields
\[
\chi^2(\mathbb Q_{u_1^{\mathrm{pt}}}\Vert\mathbb Q_{1/2})
\le e^{1/8}-1<1.
\]
For the numerical bound, use
\[
(e^{1/2})^2=e<3<4,
\qquad
e^{1/8}\le e^{1/2}<2.
\]
Consequently,
\[
\operatorname{TV}
 \bigl(\mathbb Q_{u_1^{\mathrm{pt}}},\mathbb Q_{1/2}\bigr)
\le\frac12.
\]
The testing inequality in
\cref{obj:lem:randomized-kernel-tv-comparison}, applied with the identity kernel, now gives, for every measurable
\(\mathscr A\subseteq\mathcal O_d^n\),
\[
\mathbb Q_{u_1^{\mathrm{pt}}}(\mathscr A^c)
+\mathbb Q_{u_0^{\mathrm{pt}}}(\mathscr A)
\ge
1-\operatorname{TV}
 \bigl(\mathbb Q_{u_1^{\mathrm{pt}}},
       \mathbb Q_{u_0^{\mathrm{pt}}}\bigr)
\ge\frac12.
\]
% lean: CausalSmith.Stat.MarRareqLogfrontier.oneCellFamilyLaw_obs_chi_step
% lean: CausalSmith.Stat.MarRareqLogfrontier.complete_arrival_exp_eighth_lt_two
% lean: CausalSmith.Stat.MarRareqLogfrontier.complete_arrival_product_testing_half
% lean: CausalSmith.Stat.MarRareqLogfrontier.oneCell_product_testing_half

\item \textbf{From testing to squared risk.}
Fix any \(\mathsf T\in\mathcal T_n\), and define its barycenter by
\[
b_{\mathsf T}(\boldsymbol o)
:=
\int_{[-1,1]}h\,\mathsf T(\boldsymbol o,dh)
\quad(\boldsymbol o\in\mathcal O_d^n).
\]
By \cref{obj:lem:randomized-kernel-barycenter}, this is a measurable function taking values in \([-1,1]\), and, for \(j\in\{0,1\}\),
\[
E_{\mathbb Q_{u_j^{\mathrm{pt}}}}
 \left[(b_{\mathsf T}-u_j^{\mathrm{pt}})^2\right]
\le
E_{P_{u_j^{\mathrm{pt}}}}
 \left[(\mathsf T-u_j^{\mathrm{pt}})^2\right].
\]
Define the sample events
\[
\mathscr E_j:=
\left\{\boldsymbol o:
 |b_{\mathsf T}(\boldsymbol o)-u_j^{\mathrm{pt}}|
 \ge\frac{\Delta_{\mathrm{pt}}}{2}\right\}
\quad(j=0,1),
\]
\[
\mathscr A_{\mathrm{pt}}:=
\left\{\boldsymbol o:
 |b_{\mathsf T}(\boldsymbol o)-u_1^{\mathrm{pt}}|
 \le
 |b_{\mathsf T}(\boldsymbol o)-u_0^{\mathrm{pt}}|\right\}.
\]
The triangle inequality gives, at every sample,
\[
\Delta_{\mathrm{pt}}
\le
|b_{\mathsf T}-u_0^{\mathrm{pt}}|
+
|b_{\mathsf T}-u_1^{\mathrm{pt}}|.
\]
Thus
\[
\mathscr A_{\mathrm{pt}}\subseteq\mathscr E_0,
\qquad
\mathscr A_{\mathrm{pt}}^c\subseteq\mathscr E_1,
\qquad
\mathbb Q_{u_0^{\mathrm{pt}}}(\mathscr E_0)
+\mathbb Q_{u_1^{\mathrm{pt}}}(\mathscr E_1)
\ge\frac12.
\]
Since squared error is at least \(\Delta_{\mathrm{pt}}^2/4\) on the corresponding error event,
\[
\begin{aligned}
\max_{j=0,1}
E_{P_{u_j^{\mathrm{pt}}}}
 \left[(\mathsf T-u_j^{\mathrm{pt}})^2\right]
&\ge
\frac{\Delta_{\mathrm{pt}}^2}{4}
\max_{j=0,1}
\mathbb Q_{u_j^{\mathrm{pt}}}(\mathscr E_j)\\
&\ge\frac{\Delta_{\mathrm{pt}}^2}{16}.
\end{aligned}
\]

For the scale of this bound, first suppose
\[
\frac14\le\frac1{4\sqrt N}.
\]
Then \(\Delta_{\mathrm{pt}}=1/4\), and
\[
\Delta_{\mathrm{pt}}^2=\frac1{16}
\ge\frac1{16}\min\{1,N^{-1}\}.
\]
In the complementary case,
\(\Delta_{\mathrm{pt}}=1/(4\sqrt N)\), so
\[
\Delta_{\mathrm{pt}}^2=\frac1{16N}
\ge\frac1{16}\min\{1,N^{-1}\}.
\]
Both cases therefore give
\[
\frac1{256}\min\{1,N^{-1}\}
\le
\max_{j=0,1}
E_{P_{u_j^{\mathrm{pt}}}}
 \left[(\mathsf T-u_j^{\mathrm{pt}})^2\right].
\]

Both testing laws belong to the model. The estimator class is nonempty, as witnessed by the constant-zero estimator. For every full-data law,
\[
\tau(P)=P(Y(1)=1)-P(Y(0)=1)\in[-1,1],
\]
and hence every estimator in \(\mathcal T_n\) has risk between zero and four:
\[
0\le E_P[(\mathsf T-\tau(P))^2]\le4.
\]
Thus the model supremum is a finite real number and dominates the displayed two-point maximum. Taking the infimum over estimators, as in
\cref{obj:def:minimax-risk}, proves
\[
\mathfrak R_{n,d,q}
\ge\frac1{256}\min\{1,N^{-1}\}.
\]
% lean: CausalSmith.Stat.MarRareqLogfrontier.complete_arrival_two_point_decision_reduction
% lean: CausalSmith.Stat.MarRareqLogfrontier.oneCellStep_sq_lower
% lean: CausalSmith.Stat.MarRareqLogfrontier.oneCell_two_point_risk_floor
% lean: CausalSmith.Stat.MarRareqLogfrontier.unrestricted_squaredRisk_bounds
% lean: CausalSmith.Stat.MarRareqLogfrontier.oneCell_minimax_risk_floors

\item \textbf{Construction and scale of the interval grid.}
Fix \(\alpha\in(0,1)\), and define
\[
\beta_\alpha:=1-\alpha,
\qquad
\varepsilon_\alpha:=\frac{\beta_\alpha}{8}.
\]
Choose an integer
\[
M\ge\frac8{\beta_\alpha}.
\]
Then \(M\ge1\), and \(0<\varepsilon_\alpha\le1\). Define
\[
w_{\mathrm{grid}}
:=
\min\left\{\frac14,
           \frac{\varepsilon_\alpha}{4\sqrt N}\right\},
\qquad
\Delta_{\mathrm{grid}}:=\frac{w_{\mathrm{grid}}}{M},
\qquad
u_j^{\mathrm{grid}}
:=
\frac12+j\Delta_{\mathrm{grid}}
\quad(j\in\{0,\ldots,M\}).
\]
These quantities satisfy
\[
0<w_{\mathrm{grid}}\le\frac14,
\qquad
\Delta_{\mathrm{grid}}>0,
\qquad
2Nw_{\mathrm{grid}}^2
\le
2N\left(\frac{\varepsilon_\alpha}{4\sqrt N}\right)^2
=
\frac{\varepsilon_\alpha^2}{8}.
\]
For every grid index,
\[
0\le j\Delta_{\mathrm{grid}}
=\frac{j}{M}w_{\mathrm{grid}}
\le w_{\mathrm{grid}},
\qquad
\frac12\le u_j^{\mathrm{grid}}\le\frac34.
\]
Positive spacing makes the grid points pairwise distinct.

We also need
\[
\frac{\varepsilon_\alpha}{4}
\min\{1,N^{-1/2}\}
\le w_{\mathrm{grid}}.
\]
If \(1/4\le\varepsilon_\alpha/(4\sqrt N)\), then
\[
w_{\mathrm{grid}}=\frac14,
\qquad
\frac{\varepsilon_\alpha}{4}
\min\{1,N^{-1/2}\}
\le\frac{\varepsilon_\alpha}{4}\le\frac14.
\]
In the complementary case,
\[
w_{\mathrm{grid}}
=\frac{\varepsilon_\alpha}{4\sqrt N},
\qquad
\frac{\varepsilon_\alpha}{4}
\min\{1,N^{-1/2}\}
\le\frac{\varepsilon_\alpha}{4\sqrt N}.
\]
This proves the required scale inequality in both cases.
% lean: CausalSmith.Stat.MarRareqLogfrontier.oneCell_interval_minimax_floor

\item \textbf{Coverage and comparison of interval-output laws.}
Let \(\mathsf I\in\mathcal C_{n,d,q,\alpha}\). Define its endpoint-output laws by
\[
\mu_{\mathrm{ref}}
:=
\mathbb Q_{1/2}\mathsf I,
\qquad
\mu_j:=\mathbb Q_{u_j^{\mathrm{grid}}}\mathsf I
\quad(j=0,\ldots,M),
\]
meaning explicitly that, for every measurable endpoint set \(\mathscr A_{\mathrm{end}}\),
\[
\mu_{\mathrm{ref}}(\mathscr A_{\mathrm{end}})
=
\int\mathsf I(\boldsymbol o,\mathscr A_{\mathrm{end}})\,
     \mathbb Q_{1/2}(d\boldsymbol o),
\qquad
\mu_j(\mathscr A_{\mathrm{end}})
=
\int\mathsf I(\boldsymbol o,\mathscr A_{\mathrm{end}})\,
     \mathbb Q_{u_j^{\mathrm{grid}}}(d\boldsymbol o).
\]
Write the returned endpoint pair as
\[
(l,u_{\mathrm{end}}),\qquad -1\le l\le u_{\mathrm{end}}\le1.
\]
Define the coverage events
\[
B_j:=
\{(l,u_{\mathrm{end}}):-1\le l\le u_{\mathrm{end}}\le1,\ 
          l\le u_j^{\mathrm{grid}}\le u_{\mathrm{end}}\}.
\]
All grid laws belong to the model and have the corresponding grid means as their average treatment effects. Honesty therefore gives
\[
\mu_j(B_j)\ge\beta_\alpha
\qquad(j=0,\ldots,M).
\]

For the comparison with the reference sample law, define
\[
\xi_j:=2N(j\Delta_{\mathrm{grid}})^2
\quad(j=0,\ldots,M).
\]
The width budget implies
\[
0\le\xi_j
\le2Nw_{\mathrm{grid}}^2
\le\frac{\varepsilon_\alpha^2}{8}
\le\frac18<1.
\]
The one-record identity and product calculation give
\[
\chi^2\!\left(
 \mathbb Q_{u_j^{\mathrm{grid}}}\Vert\mathbb Q_{1/2}
 \right)
\le e^{\xi_j}-1.
\]
For \(0\le\xi_j<1\), use the exponential bound
\(e^{\xi_j}\le(1-\xi_j)^{-1}\). Furthermore,
\[
\xi_j(1+\varepsilon_\alpha^2)
\le2\xi_j
\le\frac{\varepsilon_\alpha^2}{4}
\le\varepsilon_\alpha^2.
\]
Since \(1-\xi_j>0\), this implies
\[
e^{\xi_j}-1
\le\frac1{1-\xi_j}-1
=\frac{\xi_j}{1-\xi_j}
\le\varepsilon_\alpha^2.
\]
Consequently,
\[
\operatorname{TV}
 \bigl(\mathbb Q_{1/2},
       \mathbb Q_{u_j^{\mathrm{grid}}}\bigr)
\le
\frac12
\sqrt{\chi^2\!\left(
 \mathbb Q_{u_j^{\mathrm{grid}}}\Vert\mathbb Q_{1/2}
 \right)}
\le\frac{\varepsilon_\alpha}{2}
\le\varepsilon_\alpha.
\]
Applying \cref{obj:lem:randomized-kernel-tv-comparison} to the common interval kernel yields
\[
\operatorname{TV}(\mu_{\mathrm{ref}},\mu_j)
\le\varepsilon_\alpha.
\]
Thus, by the defining event bound for total variation,
\[
\mu_{\mathrm{ref}}(B_j)
\ge\mu_j(B_j)-\varepsilon_\alpha
\ge\beta_\alpha-\varepsilon_\alpha
=\frac{7\beta_\alpha}{8}.
\]
% lean: CausalSmith.Stat.MarRareqLogfrontier.oneCell_product_tv_le
% lean: Causalean.Stat.tvDist_bind_le
% lean: CausalSmith.Stat.MarRareqLogfrontier.oneCell_interval_minimax_floor

\item \textbf{Counting covered grid points.}
For every ordered endpoint pair, define
\[
C_{\mathrm{grid}}(l,u_{\mathrm{end}})
:=
\sum_{j=0}^{M}
\mathbf1\{l\le u_j^{\mathrm{grid}}\le u_{\mathrm{end}}\}
\quad(-1\le l\le u_{\mathrm{end}}\le1).
\]
We claim the pointwise inequality
\[
\Delta_{\mathrm{grid}}
\bigl(C_{\mathrm{grid}}(l,u_{\mathrm{end}})-1\bigr)
\le u_{\mathrm{end}}-l.
\]
If the covered index set is empty, its left-hand side equals
\(-\Delta_{\mathrm{grid}}\le0\le u_{\mathrm{end}}-l\). Otherwise define its extreme indices by
\[
j_-:=\min\{j\in\{0,\ldots,M\}:
                 l\le u_j^{\mathrm{grid}}\le u_{\mathrm{end}}\},
\qquad
j_+:=\max\{j\in\{0,\ldots,M\}:
                 l\le u_j^{\mathrm{grid}}\le u_{\mathrm{end}}\}.
\]
Every covered index lies between \(j_-\) and \(j_+\), so
\[
C_{\mathrm{grid}}(l,u_{\mathrm{end}})\le j_+-j_-+1.
\]
The two extreme grid points lie in the interval, and hence
\[
\Delta_{\mathrm{grid}}
 \bigl(C_{\mathrm{grid}}(l,u_{\mathrm{end}})-1\bigr)
\le
(j_+-j_-)\Delta_{\mathrm{grid}}
=
u_{j_+}^{\mathrm{grid}}-u_{j_-}^{\mathrm{grid}}
\le u_{\mathrm{end}}-l.
\]
This proves the claim, including the empty and singleton cases.

Summing coverage probabilities under the reference output law gives
\[
\int C_{\mathrm{grid}}(l,u_{\mathrm{end}})\,
     \mu_{\mathrm{ref}}(d(l,u_{\mathrm{end}}))
=
\sum_{j=0}^{M}\mu_{\mathrm{ref}}(B_j)
\ge(M+1)(\beta_\alpha-\varepsilon_\alpha).
\]
Integrating the pointwise length bound therefore yields
\[
\int(u_{\mathrm{end}}-l)\,\mu_{\mathrm{ref}}(d(l,u_{\mathrm{end}}))
\ge
\Delta_{\mathrm{grid}}
\left((M+1)\frac{7\beta_\alpha}{8}-1\right).
\]
The choice of \(M\) gives \(M\beta_\alpha\ge8\), and
\[
(M+1)\frac{7\beta_\alpha}{8}-1
-\frac{3\beta_\alpha M}{4}
=
\frac{M\beta_\alpha}{8}
+\frac{7\beta_\alpha}{8}-1
\ge0.
\]
Since \(M\Delta_{\mathrm{grid}}=w_{\mathrm{grid}}\), we conclude
\[
\int(u_{\mathrm{end}}-l)\,\mu_{\mathrm{ref}}(d(l,u_{\mathrm{end}}))
\ge\frac{3\beta_\alpha}{4}\,w_{\mathrm{grid}}.
\]
The endpoint length is measurable and bounded by two. Integrating first over the sample and then over the kernel therefore equals integration under its output law, giving
\[
E_{P_{1/2}}[u_{\mathrm{end}}-l]
=
\int(u_{\mathrm{end}}-l)\,\mu_{\mathrm{ref}}(d(l,u_{\mathrm{end}}))
\ge\frac{3\beta_\alpha}{4}\,w_{\mathrm{grid}}.
\]
All randomization is included in this expectation.
% lean: Causalean.Stat.intervalLength_ge_gridCount
% lean: Causalean.Stat.integral_gridCount
% lean: Causalean.Stat.honest_interval_length_lower_of_grid_tv_positive
% lean: CausalSmith.Stat.MarRareqLogfrontier.intervalRisk_bind_lintegral
% lean: CausalSmith.Stat.MarRareqLogfrontier.oneCell_interval_minimax_floor

\item \textbf{Minimax length and constant calibration.}
The honest interval class is nonempty. Indeed, for every full-data law,
\[
\tau(P)=P(Y(1)=1)-P(Y(0)=1)\in[-1,1],
\]
so the constant procedure returning \([-1,1]\) has coverage one. Every admissible interval procedure also satisfies, for every full-data law,
\[
0\le E_P[u_{\mathrm{end}}-l]\le2,
\]
because \(0\le u_{\mathrm{end}}-l\le2\) for every returned endpoint pair. Thus the real-valued model suprema are finite. Since \(P_{1/2}\) belongs to the model, the preceding bound implies
\[
\sup_{P\in\mathcal M_{n,d,q}}E_P[u_{\mathrm{end}}-l]
\ge E_{P_{1/2}}[u_{\mathrm{end}}-l]
\ge\frac{3\beta_\alpha}{4}\,w_{\mathrm{grid}}.
\]
Taking the infimum over honest procedures in
\cref{obj:def:interval-length-risk} gives
\[
\mathfrak L_{n,d,q,\alpha}
\ge\frac{3\beta_\alpha}{4}\,w_{\mathrm{grid}}.
\]

The width inequality and
\(\varepsilon_\alpha=\beta_\alpha/8\) give the required calibration:
\[
\begin{aligned}
\underline c_\alpha\min\{1,N^{-1/2}\}
&=
\frac{3\beta_\alpha^2}{512}
\min\{1,N^{-1/2}\}\\
&\le
\frac{3\beta_\alpha\varepsilon_\alpha}{16}
\min\{1,N^{-1/2}\}\\
&\le
\frac{3\beta_\alpha}{4}\,w_{\mathrm{grid}}.
\end{aligned}
\]
Finally, \(u_0^{\mathrm{grid}}=1/2\). Hence every honest procedure satisfies
\[
\underline c_\alpha\min\{1,N^{-1/2}\}
\le
E_{P_{u_0^{\mathrm{grid}}}}[u_{\mathrm{end}}-l]
\le
\max_{j=0,\ldots,M}E_{P_{u_j^{\mathrm{grid}}}}[u_{\mathrm{end}}-l],
\]
and the minimax length satisfies the same lower bound. These are grid points from the family fixed in the first step, completing all assertions.
% lean: CausalSmith.Stat.MarRareqLogfrontier.ate_mem_unit_interval
% lean: CausalSmith.Stat.MarRareqLogfrontier.intervalRisk_bounds
% lean: CausalSmith.Stat.MarRareqLogfrontier.oneCell_interval_minimax_floor
% lean: CausalSmith.Stat.MarRareqLogfrontier.one_cell_testing_family
\end{enumerate}
\end{proof}
